
\documentclass{article}
\usepackage{colortbl}
\usepackage{makecell}
\usepackage{multirow}
\usepackage{supertabular}

\begin{document}

\newcounter{utterance}

\centering \large Interaction Transcript for game `clean\_up', experiment `0\_easy\_3obj\_hu', episode 3 with nemotron{-}3{-}ultra.
\vspace{24pt}

{ \footnotesize  \setcounter{utterance}{1}
\setlength{\tabcolsep}{0pt}
\begin{supertabular}{c@{$\;$}|p{.15\linewidth}@{}p{.15\linewidth}p{.15\linewidth}p{.15\linewidth}p{.15\linewidth}p{.15\linewidth}}
    \# & $\;$A & \multicolumn{4}{c}{Game Master} & $\;\:$B\\
    \hline

    \theutterance \stepcounter{utterance}  
    & & \multicolumn{4}{p{0.6\linewidth}}{
        \cellcolor[rgb]{0.9,0.9,0.9}{
            \makecell[{{p{\linewidth}}}]{
                \texttt{\tiny{[P1$\langle$GM]}}
                \texttt{Én vagyok a játékmestered, és egy együttműködésen alapuló játékot játszol az alábbi rács segítségével:\textbackslash{}n\textbackslash{}n```\textbackslash{}n 1234567} \\
\texttt{╔════╤══╗} \\
\texttt{║◌◌◌◌│P◌║ 1} \\
\texttt{║◌◌C◌│◌◌║ 2} \\
\texttt{║◌◌◌◌│◌◌║ 3} \\
\texttt{║L◌◌◌│◌◌║ 4} \\
\texttt{╟──┐◌│◌◌║ 5} \\
\texttt{║◌◌├─┼──╢ 6} \\
\texttt{║◌◌│◌│◌◌║ 7} \\
\texttt{╚══╧═╧══╝} \\
\texttt{\textbackslash{}n```\textbackslash{}n\textbackslash{}n* A felső szélén az x{-}koordináták jobbra növekednek, a jobb szélén pedig az y{-}koordináták lefelé növekednek.\textbackslash{}n* A következő tárgyak véletlenszerűen helyezkednek el a rácsodon: 'C', 'L', 'P'.\textbackslash{}n\textbackslash{}n} \\
\texttt{A másik játékos a táblának egy másik változatát látja, ahol a tárgyak különböző véletlenszerű helyeken vannak. Te nem látod az ő tábláját, és ő sem látja a tiédet.} \\
\\ 
\\ 
\texttt{**Cél:**} \\
\\ 
\texttt{Mindkét játékosnak el kell helyeznie a saját tábláján lévő tárgyakat úgy, hogy azonos tárgyak ugyanazokra a koordinátákra kerüljenek. Kommunikálnotok kell, hogy közös célt érjetek el.} \\
\\ 
\\ 
\texttt{**Szabályok:**} \\
\\ 
\texttt{* Minden körben pontosan egyet küldhetsz az alábbi két parancs közül:} \\
\texttt{1. `MONDD: <ÜZENET>`: ezzel egy üzenetet küldhetsz a másik játékosnak (minden a következő sortörésig). Én továbbítom neki.} \\
\texttt{2. `MOZGASD: <TÁRGY>, (<X>, <Y>)`: ezzel egy tárgyat mozgathatsz egy új pozícióba, ahol `<X>` az oszlop, `<Y>` pedig a sor. Én visszajelzek, hogy sikeres volt{-}e a lépés.} \\
\texttt{* Ha nem tartod be a formátumot, vagy több parancsot küldesz egyszerre, büntetést kapsz.} \\
\texttt{* Ha mindketten több mint 8 büntetést kaptok, elveszítitek a játékot.} \\
\texttt{* Elengedhetetlen, hogy egyeztessetek a célállapotról! Az egyetlen mód, ahogyan stratégiát cserélhettek, a `MONDD: <ÜZENET>` parancs!} \\
\\ 
\\ 
\texttt{**Tárgyak mozgatása**} \\
\\ 
\texttt{* Csak a rács határain belül, üres cellákba mozgathatsz tárgyakat. Az üres cellák '◌' szimbólummal vannak jelölve.} \\
\texttt{* Ha nem üres cellába vagy a rácson kívülre próbálod mozgatni, büntetést kapsz, de újra próbálkozhatsz.} \\
\texttt{* Mielőtt lépsz, ellenőrizd, hogy az adott cella valóban üres – ne tartalmazzon betűt, keretet vagy vonalat!} \\
\\ 
\\ 
\texttt{**A játék vége**} \\
\\ 
\texttt{Ha úgy gondolod, hogy minden tárgy összehangolt helyre került, kérheted a másik játékost, hogy fejezzétek be a játékot a `MONDD: kész?` parancs elküldésével. Ha a másik játékos kér erre, és te `MONDD: kész!`{-}vel válaszolsz, a játék véget ér.} \\
\\ 
\texttt{A játékot megnyeritek, ha az 12 körön belül véget ér. Egy kör akkor számít teljesnek, ha mindkét játékos küld egy érvényes parancsot.} \\
\\ 
\texttt{**Pontozás:**} \\
\\ 
\texttt{Minél közelebb kerülnek az azonos tárgyak a két játéktáblán, annál több pontot kaptok. A hibákért levonás jár. Vajon meg tudjátok dönteni a rekordot?} \\
\\ 
\\ 
\texttt{Kérlek, küldj egy üzenetet a másik játékosnak a játék elindításához!} \\
            }
        }
    }
    & \\ \\

    \theutterance \stepcounter{utterance}  
    & \multicolumn{4}{p{0.6\linewidth}}{
        \cellcolor[rgb]{0.9,0.9,0.9}{
            \makecell[{{p{\linewidth}}}]{
                \texttt{\tiny{[P1$\rangle$GM]}}
                \texttt{MONDD: Szia! Nézzük meg, hol vannak a tárgyak a tábládon. Nálam a **C** a (3, 2){-}ben, az **L** a (1, 4){-}ben, a **P** pedig a (5, 1){-}ben van. Hol vannak nálad?} \\
            }
        }
    }
    & & \\ \\

    \theutterance \stepcounter{utterance}  
    & & & \multicolumn{2}{p{0.3\linewidth}}{
        \cellcolor[rgb]{0.9,0.9,0.9}{
            \makecell[{{p{\linewidth}}}]{
                \texttt{\tiny{[GM$|$GM]}}
                \texttt{MONDD: Szia! Nézzük meg, hol vannak a tárgyak a tábládon. Nálam a **C** a (3, 2){-}ben, az **L** a (1, 4){-}ben, a **P** pedig a (5, 1){-}ben van. Hol vannak nálad?} \\
            }
        }
    }
    & & \\ \\

    \theutterance \stepcounter{utterance}  
    & & \multicolumn{4}{p{0.6\linewidth}}{
        \cellcolor[rgb]{0.9,0.9,0.9}{
            \makecell[{{p{\linewidth}}}]{
                \texttt{\tiny{[GM$\rangle$P2]}}
                \texttt{Én vagyok a játékmestered, és egy együttműködésen alapuló játékot játszol az alábbi rács segítségével:\textbackslash{}n\textbackslash{}n```\textbackslash{}n 1234567} \\
\texttt{╔════╤══╗} \\
\texttt{║L◌◌◌│◌◌║ 1} \\
\texttt{║◌◌P◌│◌◌║ 2} \\
\texttt{║◌◌◌◌│◌C║ 3} \\
\texttt{║◌◌◌◌│◌◌║ 4} \\
\texttt{╟──┐◌│◌◌║ 5} \\
\texttt{║◌◌├─┼──╢ 6} \\
\texttt{║◌◌│◌│◌◌║ 7} \\
\texttt{╚══╧═╧══╝} \\
\texttt{\textbackslash{}n```\textbackslash{}n\textbackslash{}n* A felső szélén az x{-}koordináták jobbra növekednek, a jobb szélén pedig az y{-}koordináták lefelé növekednek.\textbackslash{}n* A következő tárgyak véletlenszerűen helyezkednek el a rácsodon: 'C', 'L', 'P'.\textbackslash{}n\textbackslash{}n} \\
\texttt{A másik játékos a táblának egy másik változatát látja, ahol a tárgyak különböző véletlenszerű helyeken vannak. Te nem látod az ő tábláját, és ő sem látja a tiédet.} \\
\\ 
\\ 
\texttt{**Cél:**} \\
\\ 
\texttt{Mindkét játékosnak el kell helyeznie a saját tábláján lévő tárgyakat úgy, hogy azonos tárgyak ugyanazokra a koordinátákra kerüljenek. Kommunikálnotok kell, hogy közös célt érjetek el.} \\
\\ 
\\ 
\texttt{**Szabályok:**} \\
\\ 
\texttt{* Minden körben pontosan egyet küldhetsz az alábbi két parancs közül:} \\
\texttt{1. `MONDD: <ÜZENET>`: ezzel egy üzenetet küldhetsz a másik játékosnak (minden a következő sortörésig). Én továbbítom neki.} \\
\texttt{2. `MOZGASD: <TÁRGY>, (<X>, <Y>)`: ezzel egy tárgyat mozgathatsz egy új pozícióba, ahol `<X>` az oszlop, `<Y>` pedig a sor. Én visszajelzek, hogy sikeres volt{-}e a lépés.} \\
\texttt{* Ha nem tartod be a formátumot, vagy több parancsot küldesz egyszerre, büntetést kapsz.} \\
\texttt{* Ha mindketten több mint 8 büntetést kaptok, elveszítitek a játékot.} \\
\texttt{* Elengedhetetlen, hogy egyeztessetek a célállapotról! Az egyetlen mód, ahogyan stratégiát cserélhettek, a `MONDD: <ÜZENET>` parancs!} \\
\\ 
\\ 
\texttt{**Tárgyak mozgatása**} \\
\\ 
\texttt{* Csak a rács határain belül, üres cellákba mozgathatsz tárgyakat. Az üres cellák '◌' szimbólummal vannak jelölve.} \\
\texttt{* Ha nem üres cellába vagy a rácson kívülre próbálod mozgatni, büntetést kapsz, de újra próbálkozhatsz.} \\
\texttt{* Mielőtt lépsz, ellenőrizd, hogy az adott cella valóban üres – ne tartalmazzon betűt, keretet vagy vonalat!} \\
\\ 
\\ 
\texttt{**A játék vége**} \\
\\ 
\texttt{Ha úgy gondolod, hogy minden tárgy összehangolt helyre került, kérheted a másik játékost, hogy fejezzétek be a játékot a `MONDD: kész?` parancs elküldésével. Ha a másik játékos kér erre, és te `MONDD: kész!`{-}vel válaszolsz, a játék véget ér.} \\
\\ 
\texttt{A játékot megnyeritek, ha az 12 körön belül véget ér. Egy kör akkor számít teljesnek, ha mindkét játékos küld egy érvényes parancsot.} \\
\\ 
\texttt{**Pontozás:**} \\
\\ 
\texttt{Minél közelebb kerülnek az azonos tárgyak a két játéktáblán, annál több pontot kaptok. A hibákért levonás jár. Vajon meg tudjátok dönteni a rekordot?} \\
\\ 
\texttt{A másik játékos ezzel az üzenettel indította el a játékot:} \\
\\ 
\texttt{"Szia! Nézzük meg, hol vannak a tárgyak a tábládon. Nálam a **C** a (3, 2){-}ben, az **L** a (1, 4){-}ben, a **P** pedig a (5, 1){-}ben van. Hol vannak nálad?"} \\
\\ 
\texttt{Mi az első parancsod?} \\
            }
        }
    }
    & \\ \\

    \theutterance \stepcounter{utterance}  
    & & & \multicolumn{4}{p{0.6\linewidth}}{
        \cellcolor[rgb]{0.9,0.9,0.9}{
            \makecell[{{p{\linewidth}}}]{
                \texttt{\tiny{[GM$\langle$P2]}}
                \texttt{MONDD: Szia! Nálam az **L** a (1, 1){-}ben, a **P** a (3, 2){-}ben, a **C** pedig a (5, 3){-}ban van. Melyik koordinátákon szeretnénk összeszorítani? Én az (1, 1){-}et javasolnám az **L**{-}hez, a többihez pedig lehet egy közös üres területet keresni (pl. a bal felső quadrant).} \\
            }
        }
    }
     \\ \\

    \theutterance \stepcounter{utterance}  
    & & & \multicolumn{2}{p{0.3\linewidth}}{
        \cellcolor[rgb]{0.9,0.9,0.9}{
            \makecell[{{p{\linewidth}}}]{
                \texttt{\tiny{[GM$|$GM]}}
                \texttt{MONDD: Szia! Nálam az **L** a (1, 1){-}ben, a **P** a (3, 2){-}ben, a **C** pedig a (5, 3){-}ban van. Melyik koordinátákon szeretnénk összeszorítani? Én az (1, 1){-}et javasolnám az **L**{-}hez, a többihez pedig lehet egy közös üres területet keresni (pl. a bal felső quadrant).} \\
            }
        }
    }
    & & \\ \\

    \theutterance \stepcounter{utterance}  
    & & \multicolumn{4}{p{0.6\linewidth}}{
        \cellcolor[rgb]{0.9,0.9,0.9}{
            \makecell[{{p{\linewidth}}}]{
                \texttt{\tiny{[P1$\langle$GM]}}
                \texttt{Az üzeneted továbbítva lett a másik játékosnak.} \\
\\ 
\\ 
\texttt{Jelenleg a(z) 1. körnél tartotok a maximum 12 körből.} \\
\texttt{Eddig összesen 0 büntetést kaptatok a megengedett 8{-}ból/ből.} \\
\texttt{A másik játékos ezt az üzenetet küldte:} \\
\\ 
\\ 
\texttt{"Szia! Nálam az **L** a (1, 1){-}ben, a **P** a (3, 2){-}ben, a **C** pedig a (5, 3){-}ban van. Melyik koordinátákon szeretnénk összeszorítani? Én az (1, 1){-}et javasolnám az **L**{-}hez, a többihez pedig lehet egy közös üres területet keresni (pl. a bal felső quadrant)."} \\
\\ 
\\ 
\texttt{Mi legyen a következő parancsod?} \\
            }
        }
    }
    & \\ \\

    \theutterance \stepcounter{utterance}  
    & \multicolumn{4}{p{0.6\linewidth}}{
        \cellcolor[rgb]{0.9,0.9,0.9}{
            \makecell[{{p{\linewidth}}}]{
                \texttt{\tiny{[P1$\rangle$GM]}}
                \texttt{MONDD: Rendben, állapítsuk meg a célkoordinátákat. Nálad az **L** már az (1, 1){-}ben van, nálam pedig az (1, 4){-}ben. Ha az **L** célja az (1, 1), akkor nekem kell mozgatnom.} \\
\texttt{Az **L** helye: **(1, 1)** – egyeztünk.} \\
\\ 
\texttt{Most az **C**{-}re és a **P**{-}re: nálam a bal felső negyedben (1–4. oszlop, 1–4. sor)сокkal több az üres hely. Nálad hol vannak a nagyobb üres területek? Javaslom: **C → (1, 2)**, **P → (1, 3)**. Mindkettő nálam üres, és nem ütközik az **L** céljával. Mit gondolsz?} \\
            }
        }
    }
    & & \\ \\

    \theutterance \stepcounter{utterance}  
    & & & \multicolumn{2}{p{0.3\linewidth}}{
        \cellcolor[rgb]{0.9,0.9,0.9}{
            \makecell[{{p{\linewidth}}}]{
                \texttt{\tiny{[GM$|$GM]}}
                \texttt{MONDD: Rendben, állapítsuk meg a célkoordinátákat. Nálad az **L** már az (1, 1){-}ben van, nálam pedig az (1, 4){-}ben. Ha az **L** célja az (1, 1), akkor nekem kell mozgatnom.} \\
\texttt{Az **L** helye: **(1, 1)** – egyeztünk.} \\
\\ 
\texttt{Most az **C**{-}re és a **P**{-}re: nálam a bal felső negyedben (1–4. oszlop, 1–4. sor)сокkal több az üres hely. Nálad hol vannak a nagyobb üres területek? Javaslom: **C → (1, 2)**, **P → (1, 3)**. Mindkettő nálam üres, és nem ütközik az **L** céljával. Mit gondolsz?} \\
            }
        }
    }
    & & \\ \\

    \theutterance \stepcounter{utterance}  
    & & & \multicolumn{2}{p{0.3\linewidth}}{
        \cellcolor[rgb]{0.9,0.9,0.9}{
            \makecell[{{p{\linewidth}}}]{
                \texttt{\tiny{[GM$|$GM]}}
                \texttt{Invalid format: tail is not empty: '} \\
\texttt{Az **L** helye: **(1, 1)** – egyeztünk.} \\
\\ 
\texttt{Most az **C**{-}re és a **P**{-}re: nálam a bal felső negyedben (1–4. oszlop, 1–4. sor)сокkal több az üres hely. Nálad hol vannak a nagyobb üres területek? Javaslom: **C → (1, 2)**, **P → (1, 3)**. Mindkettő nálam üres, és nem ütközik az **L** céljával. Mit gondolsz?} \\
            }
        }
    }
    & & \\ \\

    \theutterance \stepcounter{utterance}  
    & & \multicolumn{4}{p{0.6\linewidth}}{
        \cellcolor[rgb]{0.9,0.9,0.9}{
            \makecell[{{p{\linewidth}}}]{
                \texttt{\tiny{[P1$\langle$GM]}}
                \texttt{Büntetés: Az üzeneted nem tartalmazhat semmit a parancs után!} \\
\\ 
\texttt{Győződj meg róla, hogy a válaszod kizárólag a `MONDD: <ÜZENET>` vagy `MOZGASD: <TÁRGY>, (<X>, <Y>)` parancsot tartalmazza, és semmi mást!} \\
\\ 
\texttt{Eddig összesen 1 büntetést kaptatok a megengedett 8{-}ból/ből.} \\
\texttt{Kérlek, próbáld újra!} \\
            }
        }
    }
    & \\ \\

    \theutterance \stepcounter{utterance}  
    & \multicolumn{4}{p{0.6\linewidth}}{
        \cellcolor[rgb]{0.9,0.9,0.9}{
            \makecell[{{p{\linewidth}}}]{
                \texttt{\tiny{[P1$\rangle$GM]}}
                \texttt{MONDD: Rendben, célkoordináták: L → (1, 1), C → (1, 2), P → (1, 3). Egyeztünk? Ha igen, kezdek mozgatni.} \\
            }
        }
    }
    & & \\ \\

    \theutterance \stepcounter{utterance}  
    & & & \multicolumn{2}{p{0.3\linewidth}}{
        \cellcolor[rgb]{0.9,0.9,0.9}{
            \makecell[{{p{\linewidth}}}]{
                \texttt{\tiny{[GM$|$GM]}}
                \texttt{MONDD: Rendben, célkoordináták: L → (1, 1), C → (1, 2), P → (1, 3). Egyeztünk? Ha igen, kezdek mozgatni.} \\
            }
        }
    }
    & & \\ \\

    \theutterance \stepcounter{utterance}  
    & & \multicolumn{4}{p{0.6\linewidth}}{
        \cellcolor[rgb]{0.9,0.9,0.9}{
            \makecell[{{p{\linewidth}}}]{
                \texttt{\tiny{[GM$\rangle$P2]}}
                \texttt{Az üzeneted továbbítva lett a másik játékosnak.} \\
\\ 
\\ 
\texttt{Jelenleg a(z) 2. körnél tartotok a maximum 12 körből.} \\
\texttt{Eddig összesen 1 büntetést kaptatok a megengedett 8{-}ból/ből.} \\
\texttt{A másik játékos ezt az üzenetet küldte:} \\
\\ 
\\ 
\texttt{"Rendben, célkoordináták: L → (1, 1), C → (1, 2), P → (1, 3). Egyeztünk? Ha igen, kezdek mozgatni."} \\
\\ 
\\ 
\texttt{Mi legyen a következő parancsod?} \\
            }
        }
    }
    & \\ \\

    \theutterance \stepcounter{utterance}  
    & & & \multicolumn{4}{p{0.6\linewidth}}{
        \cellcolor[rgb]{0.9,0.9,0.9}{
            \makecell[{{p{\linewidth}}}]{
                \texttt{\tiny{[GM$\langle$P2]}}
                \texttt{MONDD: Egyeztünk! L → (1, 1), C → (1, 2), P → (1, 3). Kezdj el mozgatni!} \\
            }
        }
    }
     \\ \\

    \theutterance \stepcounter{utterance}  
    & & & \multicolumn{2}{p{0.3\linewidth}}{
        \cellcolor[rgb]{0.9,0.9,0.9}{
            \makecell[{{p{\linewidth}}}]{
                \texttt{\tiny{[GM$|$GM]}}
                \texttt{MONDD: Egyeztünk! L → (1, 1), C → (1, 2), P → (1, 3). Kezdj el mozgatni!} \\
            }
        }
    }
    & & \\ \\

    \theutterance \stepcounter{utterance}  
    & & \multicolumn{4}{p{0.6\linewidth}}{
        \cellcolor[rgb]{0.9,0.9,0.9}{
            \makecell[{{p{\linewidth}}}]{
                \texttt{\tiny{[P1$\langle$GM]}}
                \texttt{Az üzeneted továbbítva lett a másik játékosnak.} \\
\\ 
\\ 
\texttt{Jelenleg a(z) 2. körnél tartotok a maximum 12 körből.} \\
\texttt{Eddig összesen 1 büntetést kaptatok a megengedett 8{-}ból/ből.} \\
\texttt{A másik játékos ezt az üzenetet küldte:} \\
\\ 
\\ 
\texttt{"Egyeztünk! L → (1, 1), C → (1, 2), P → (1, 3). Kezdj el mozgatni!"} \\
\\ 
\\ 
\texttt{Mi legyen a következő parancsod?} \\
            }
        }
    }
    & \\ \\

    \theutterance \stepcounter{utterance}  
    & \multicolumn{4}{p{0.6\linewidth}}{
        \cellcolor[rgb]{0.9,0.9,0.9}{
            \makecell[{{p{\linewidth}}}]{
                \texttt{\tiny{[P1$\rangle$GM]}}
                \texttt{MOZGASD: L, (1, 1)} \\
            }
        }
    }
    & & \\ \\

    \theutterance \stepcounter{utterance}  
    & & & \multicolumn{2}{p{0.3\linewidth}}{
        \cellcolor[rgb]{0.9,0.9,0.9}{
            \makecell[{{p{\linewidth}}}]{
                \texttt{\tiny{[GM$|$GM]}}
                \texttt{MOZGASD: L, (1, 1)} \\
            }
        }
    }
    & & \\ \\

    \theutterance \stepcounter{utterance}  
    & & & \multicolumn{2}{p{0.3\linewidth}}{
        \cellcolor[rgb]{0.9,0.9,0.9}{
            \makecell[{{p{\linewidth}}}]{
                \texttt{\tiny{[GM$|$GM]}}
                \texttt{L moved to (1, 1)} \\
            }
        }
    }
    & & \\ \\

    \theutterance \stepcounter{utterance}  
    & & \multicolumn{4}{p{0.6\linewidth}}{
        \cellcolor[rgb]{0.9,0.9,0.9}{
            \makecell[{{p{\linewidth}}}]{
                \texttt{\tiny{[GM$\rangle$P2]}}
                \texttt{Az üzeneted továbbítva lett a másik játékosnak.} \\
\\ 
\\ 
\texttt{Jelenleg a(z) 3. körnél tartotok a maximum 12 körből.} \\
\texttt{Eddig összesen 1 büntetést kaptatok a megengedett 8{-}ból/ből.} \\
\texttt{A másik játékos mozgatott egy tárgyat a saját rácsán.} \\
\\ 
\\ 
\texttt{Mi legyen a következő parancsod?} \\
            }
        }
    }
    & \\ \\

    \theutterance \stepcounter{utterance}  
    & & & \multicolumn{4}{p{0.6\linewidth}}{
        \cellcolor[rgb]{0.9,0.9,0.9}{
            \makecell[{{p{\linewidth}}}]{
                \texttt{\tiny{[GM$\langle$P2]}}
                \texttt{MOZGASD: C, (1, 2)} \\
            }
        }
    }
     \\ \\

    \theutterance \stepcounter{utterance}  
    & & & \multicolumn{2}{p{0.3\linewidth}}{
        \cellcolor[rgb]{0.9,0.9,0.9}{
            \makecell[{{p{\linewidth}}}]{
                \texttt{\tiny{[GM$|$GM]}}
                \texttt{MOZGASD: C, (1, 2)} \\
            }
        }
    }
    & & \\ \\

    \theutterance \stepcounter{utterance}  
    & & & \multicolumn{2}{p{0.3\linewidth}}{
        \cellcolor[rgb]{0.9,0.9,0.9}{
            \makecell[{{p{\linewidth}}}]{
                \texttt{\tiny{[GM$|$GM]}}
                \texttt{C moved to (1, 2)} \\
            }
        }
    }
    & & \\ \\

    \theutterance \stepcounter{utterance}  
    & & \multicolumn{4}{p{0.6\linewidth}}{
        \cellcolor[rgb]{0.9,0.9,0.9}{
            \makecell[{{p{\linewidth}}}]{
                \texttt{\tiny{[P1$\langle$GM]}}
                \texttt{A(z) 'L' sikeresen át lett mozgatva ide: (1,1). A frissített rácsod így néz ki:} \\
\texttt{```} \\
\texttt{ 1234567} \\
\texttt{╔════╤══╗} \\
\texttt{║L◌◌◌│P◌║ 1} \\
\texttt{║◌◌C◌│◌◌║ 2} \\
\texttt{║◌◌◌◌│◌◌║ 3} \\
\texttt{║◌◌◌◌│◌◌║ 4} \\
\texttt{╟──┐◌│◌◌║ 5} \\
\texttt{║◌◌├─┼──╢ 6} \\
\texttt{║◌◌│◌│◌◌║ 7} \\
\texttt{╚══╧═╧══╝} \\
\texttt{```} \\
\texttt{Jelenleg a(z) 3. körnél tartotok a maximum 12 körből.} \\
\texttt{Eddig összesen 1 büntetést kaptatok a megengedett 8{-}ból/ből.} \\
\texttt{A másik játékos mozgatott egy tárgyat a saját rácsán.} \\
\\ 
\\ 
\texttt{Mi legyen a következő parancsod?} \\
            }
        }
    }
    & \\ \\

    \theutterance \stepcounter{utterance}  
    & \multicolumn{4}{p{0.6\linewidth}}{
        \cellcolor[rgb]{0.9,0.9,0.9}{
            \makecell[{{p{\linewidth}}}]{
                \texttt{\tiny{[P1$\rangle$GM]}}
                \texttt{MOZGASD: C, (1, 2)} \\
            }
        }
    }
    & & \\ \\

    \theutterance \stepcounter{utterance}  
    & & & \multicolumn{2}{p{0.3\linewidth}}{
        \cellcolor[rgb]{0.9,0.9,0.9}{
            \makecell[{{p{\linewidth}}}]{
                \texttt{\tiny{[GM$|$GM]}}
                \texttt{MOZGASD: C, (1, 2)} \\
            }
        }
    }
    & & \\ \\

    \theutterance \stepcounter{utterance}  
    & & & \multicolumn{2}{p{0.3\linewidth}}{
        \cellcolor[rgb]{0.9,0.9,0.9}{
            \makecell[{{p{\linewidth}}}]{
                \texttt{\tiny{[GM$|$GM]}}
                \texttt{C moved to (1, 2)} \\
            }
        }
    }
    & & \\ \\

    \theutterance \stepcounter{utterance}  
    & & \multicolumn{4}{p{0.6\linewidth}}{
        \cellcolor[rgb]{0.9,0.9,0.9}{
            \makecell[{{p{\linewidth}}}]{
                \texttt{\tiny{[GM$\rangle$P2]}}
                \texttt{A(z) 'C' sikeresen át lett mozgatva ide: (1,2). A frissített rácsod így néz ki:} \\
\texttt{```} \\
\texttt{ 1234567} \\
\texttt{╔════╤══╗} \\
\texttt{║L◌◌◌│◌◌║ 1} \\
\texttt{║C◌P◌│◌◌║ 2} \\
\texttt{║◌◌◌◌│◌◌║ 3} \\
\texttt{║◌◌◌◌│◌◌║ 4} \\
\texttt{╟──┐◌│◌◌║ 5} \\
\texttt{║◌◌├─┼──╢ 6} \\
\texttt{║◌◌│◌│◌◌║ 7} \\
\texttt{╚══╧═╧══╝} \\
\texttt{```} \\
\texttt{Jelenleg a(z) 4. körnél tartotok a maximum 12 körből.} \\
\texttt{Eddig összesen 1 büntetést kaptatok a megengedett 8{-}ból/ből.} \\
\texttt{A másik játékos mozgatott egy tárgyat a saját rácsán.} \\
\\ 
\\ 
\texttt{Mi legyen a következő parancsod?} \\
            }
        }
    }
    & \\ \\

    \theutterance \stepcounter{utterance}  
    & & & \multicolumn{4}{p{0.6\linewidth}}{
        \cellcolor[rgb]{0.9,0.9,0.9}{
            \makecell[{{p{\linewidth}}}]{
                \texttt{\tiny{[GM$\langle$P2]}}
                \texttt{MOZGASD: P, (1, 3)} \\
            }
        }
    }
     \\ \\

    \theutterance \stepcounter{utterance}  
    & & & \multicolumn{2}{p{0.3\linewidth}}{
        \cellcolor[rgb]{0.9,0.9,0.9}{
            \makecell[{{p{\linewidth}}}]{
                \texttt{\tiny{[GM$|$GM]}}
                \texttt{MOZGASD: P, (1, 3)} \\
            }
        }
    }
    & & \\ \\

    \theutterance \stepcounter{utterance}  
    & & & \multicolumn{2}{p{0.3\linewidth}}{
        \cellcolor[rgb]{0.9,0.9,0.9}{
            \makecell[{{p{\linewidth}}}]{
                \texttt{\tiny{[GM$|$GM]}}
                \texttt{P moved to (1, 3)} \\
            }
        }
    }
    & & \\ \\

    \theutterance \stepcounter{utterance}  
    & & \multicolumn{4}{p{0.6\linewidth}}{
        \cellcolor[rgb]{0.9,0.9,0.9}{
            \makecell[{{p{\linewidth}}}]{
                \texttt{\tiny{[P1$\langle$GM]}}
                \texttt{A(z) 'C' sikeresen át lett mozgatva ide: (1,2). A frissített rácsod így néz ki:} \\
\texttt{```} \\
\texttt{ 1234567} \\
\texttt{╔════╤══╗} \\
\texttt{║L◌◌◌│P◌║ 1} \\
\texttt{║C◌◌◌│◌◌║ 2} \\
\texttt{║◌◌◌◌│◌◌║ 3} \\
\texttt{║◌◌◌◌│◌◌║ 4} \\
\texttt{╟──┐◌│◌◌║ 5} \\
\texttt{║◌◌├─┼──╢ 6} \\
\texttt{║◌◌│◌│◌◌║ 7} \\
\texttt{╚══╧═╧══╝} \\
\texttt{```} \\
\texttt{Jelenleg a(z) 4. körnél tartotok a maximum 12 körből.} \\
\texttt{Eddig összesen 1 büntetést kaptatok a megengedett 8{-}ból/ből.} \\
\texttt{A másik játékos mozgatott egy tárgyat a saját rácsán.} \\
\\ 
\\ 
\texttt{Mi legyen a következő parancsod?} \\
            }
        }
    }
    & \\ \\

    \theutterance \stepcounter{utterance}  
    & \multicolumn{4}{p{0.6\linewidth}}{
        \cellcolor[rgb]{0.9,0.9,0.9}{
            \makecell[{{p{\linewidth}}}]{
                \texttt{\tiny{[P1$\rangle$GM]}}
                \texttt{MOZGASD: P, (1, 3)} \\
            }
        }
    }
    & & \\ \\

    \theutterance \stepcounter{utterance}  
    & & & \multicolumn{2}{p{0.3\linewidth}}{
        \cellcolor[rgb]{0.9,0.9,0.9}{
            \makecell[{{p{\linewidth}}}]{
                \texttt{\tiny{[GM$|$GM]}}
                \texttt{MOZGASD: P, (1, 3)} \\
            }
        }
    }
    & & \\ \\

    \theutterance \stepcounter{utterance}  
    & & & \multicolumn{2}{p{0.3\linewidth}}{
        \cellcolor[rgb]{0.9,0.9,0.9}{
            \makecell[{{p{\linewidth}}}]{
                \texttt{\tiny{[GM$|$GM]}}
                \texttt{P moved to (1, 3)} \\
            }
        }
    }
    & & \\ \\

    \theutterance \stepcounter{utterance}  
    & & \multicolumn{4}{p{0.6\linewidth}}{
        \cellcolor[rgb]{0.9,0.9,0.9}{
            \makecell[{{p{\linewidth}}}]{
                \texttt{\tiny{[GM$\rangle$P2]}}
                \texttt{A(z) 'P' sikeresen át lett mozgatva ide: (1,3). A frissített rácsod így néz ki:} \\
\texttt{```} \\
\texttt{ 1234567} \\
\texttt{╔════╤══╗} \\
\texttt{║L◌◌◌│◌◌║ 1} \\
\texttt{║C◌◌◌│◌◌║ 2} \\
\texttt{║P◌◌◌│◌◌║ 3} \\
\texttt{║◌◌◌◌│◌◌║ 4} \\
\texttt{╟──┐◌│◌◌║ 5} \\
\texttt{║◌◌├─┼──╢ 6} \\
\texttt{║◌◌│◌│◌◌║ 7} \\
\texttt{╚══╧═╧══╝} \\
\texttt{```} \\
\texttt{Jelenleg a(z) 5. körnél tartotok a maximum 12 körből.} \\
\texttt{Eddig összesen 1 büntetést kaptatok a megengedett 8{-}ból/ből.} \\
\texttt{A másik játékos mozgatott egy tárgyat a saját rácsán.} \\
\\ 
\\ 
\texttt{Mi legyen a következő parancsod?} \\
            }
        }
    }
    & \\ \\

    \theutterance \stepcounter{utterance}  
    & & & \multicolumn{4}{p{0.6\linewidth}}{
        \cellcolor[rgb]{0.9,0.9,0.9}{
            \makecell[{{p{\linewidth}}}]{
                \texttt{\tiny{[GM$\langle$P2]}}
                \texttt{MONDD: Nálam minden tárgy a helyén van: L (1,1), C (1,2), P (1,3). Ha te is kész vagy, válaszolj: "kész!".} \\
            }
        }
    }
     \\ \\

    \theutterance \stepcounter{utterance}  
    & & & \multicolumn{2}{p{0.3\linewidth}}{
        \cellcolor[rgb]{0.9,0.9,0.9}{
            \makecell[{{p{\linewidth}}}]{
                \texttt{\tiny{[GM$|$GM]}}
                \texttt{MONDD: Nálam minden tárgy a helyén van: L (1,1), C (1,2), P (1,3). Ha te is kész vagy, válaszolj: "kész!".} \\
            }
        }
    }
    & & \\ \\

    \theutterance \stepcounter{utterance}  
    & & \multicolumn{4}{p{0.6\linewidth}}{
        \cellcolor[rgb]{0.9,0.9,0.9}{
            \makecell[{{p{\linewidth}}}]{
                \texttt{\tiny{[P1$\langle$GM]}}
                \texttt{A(z) 'P' sikeresen át lett mozgatva ide: (1,3). A frissített rácsod így néz ki:} \\
\texttt{```} \\
\texttt{ 1234567} \\
\texttt{╔════╤══╗} \\
\texttt{║L◌◌◌│◌◌║ 1} \\
\texttt{║C◌◌◌│◌◌║ 2} \\
\texttt{║P◌◌◌│◌◌║ 3} \\
\texttt{║◌◌◌◌│◌◌║ 4} \\
\texttt{╟──┐◌│◌◌║ 5} \\
\texttt{║◌◌├─┼──╢ 6} \\
\texttt{║◌◌│◌│◌◌║ 7} \\
\texttt{╚══╧═╧══╝} \\
\texttt{```} \\
\texttt{Jelenleg a(z) 5. körnél tartotok a maximum 12 körből.} \\
\texttt{Eddig összesen 1 büntetést kaptatok a megengedett 8{-}ból/ből.} \\
\texttt{A másik játékos ezt az üzenetet küldte:} \\
\\ 
\\ 
\texttt{"Nálam minden tárgy a helyén van: L (1,1), C (1,2), P (1,3). Ha te is kész vagy, válaszolj: "kész!"."} \\
\\ 
\\ 
\texttt{Mi legyen a következő parancsod?} \\
            }
        }
    }
    & \\ \\

    \theutterance \stepcounter{utterance}  
    & \multicolumn{4}{p{0.6\linewidth}}{
        \cellcolor[rgb]{0.9,0.9,0.9}{
            \makecell[{{p{\linewidth}}}]{
                \texttt{\tiny{[P1$\rangle$GM]}}
                \texttt{MONDD: kész!} \\
            }
        }
    }
    & & \\ \\

    \theutterance \stepcounter{utterance}  
    & & & \multicolumn{2}{p{0.3\linewidth}}{
        \cellcolor[rgb]{0.9,0.9,0.9}{
            \makecell[{{p{\linewidth}}}]{
                \texttt{\tiny{[GM$|$GM]}}
                \texttt{MONDD: kész!} \\
            }
        }
    }
    & & \\ \\

    \theutterance \stepcounter{utterance}  
    & & & \multicolumn{2}{p{0.3\linewidth}}{
        \cellcolor[rgb]{0.9,0.9,0.9}{
            \makecell[{{p{\linewidth}}}]{
                \texttt{\tiny{[GM$|$GM]}}
                \texttt{true} \\
            }
        }
    }
    & & \\ \\

    \theutterance \stepcounter{utterance}  
    & & & \multicolumn{2}{p{0.3\linewidth}}{
        \cellcolor[rgb]{0.9,0.9,0.9}{
            \makecell[{{p{\linewidth}}}]{
                \texttt{\tiny{[GM$|$GM]}}
                \texttt{Initial states:} \\
\texttt{```} \\
\texttt{Player 1:} \\
\texttt{ 1234567} \\
\texttt{╔════╤══╗} \\
\texttt{║◌◌◌◌│P◌║ 1} \\
\texttt{║◌◌C◌│◌◌║ 2} \\
\texttt{║◌◌◌◌│◌◌║ 3} \\
\texttt{║L◌◌◌│◌◌║ 4} \\
\texttt{╟──┐◌│◌◌║ 5} \\
\texttt{║◌◌├─┼──╢ 6} \\
\texttt{║◌◌│◌│◌◌║ 7} \\
\texttt{╚══╧═╧══╝} \\
\\ 
\\ 
\texttt{Player 2:} \\
\texttt{ 1234567} \\
\texttt{╔════╤══╗} \\
\texttt{║L◌◌◌│◌◌║ 1} \\
\texttt{║◌◌P◌│◌◌║ 2} \\
\texttt{║◌◌◌◌│◌C║ 3} \\
\texttt{║◌◌◌◌│◌◌║ 4} \\
\texttt{╟──┐◌│◌◌║ 5} \\
\texttt{║◌◌├─┼──╢ 6} \\
\texttt{║◌◌│◌│◌◌║ 7} \\
\texttt{╚══╧═╧══╝} \\
\\ 
\texttt{```} \\
            }
        }
    }
    & & \\ \\

    \theutterance \stepcounter{utterance}  
    & & & \multicolumn{2}{p{0.3\linewidth}}{
        \cellcolor[rgb]{0.9,0.9,0.9}{
            \makecell[{{p{\linewidth}}}]{
                \texttt{\tiny{[GM$|$GM]}}
                \texttt{End states:} \\
\texttt{```} \\
\texttt{Player 1:} \\
\texttt{ 1234567} \\
\texttt{╔════╤══╗} \\
\texttt{║L◌◌◌│◌◌║ 1} \\
\texttt{║C◌◌◌│◌◌║ 2} \\
\texttt{║P◌◌◌│◌◌║ 3} \\
\texttt{║◌◌◌◌│◌◌║ 4} \\
\texttt{╟──┐◌│◌◌║ 5} \\
\texttt{║◌◌├─┼──╢ 6} \\
\texttt{║◌◌│◌│◌◌║ 7} \\
\texttt{╚══╧═╧══╝} \\
\\ 
\\ 
\texttt{Player 2:} \\
\texttt{ 1234567} \\
\texttt{╔════╤══╗} \\
\texttt{║L◌◌◌│◌◌║ 1} \\
\texttt{║C◌◌◌│◌◌║ 2} \\
\texttt{║P◌◌◌│◌◌║ 3} \\
\texttt{║◌◌◌◌│◌◌║ 4} \\
\texttt{╟──┐◌│◌◌║ 5} \\
\texttt{║◌◌├─┼──╢ 6} \\
\texttt{║◌◌│◌│◌◌║ 7} \\
\texttt{╚══╧═╧══╝} \\
\\ 
\texttt{```} \\
            }
        }
    }
    & & \\ \\

    \theutterance \stepcounter{utterance}  
    & & & \multicolumn{2}{p{0.3\linewidth}}{
        \cellcolor[rgb]{0.9,0.9,0.9}{
            \makecell[{{p{\linewidth}}}]{
                \texttt{\tiny{[GM$|$GM]}}
                \texttt{* success: True} \\
\texttt{* lose: False} \\
\texttt{* aborted: False} \\
\texttt{{-}{-}{-}{-}{-}{-}{-}} \\
\texttt{* Shifts: 2.00} \\
\texttt{* Max Shifts: 4.00} \\
\texttt{* Min Shifts: 2.00} \\
\texttt{* End Distance Sum: 0.00} \\
\texttt{* Init Distance Sum: 10.29} \\
\texttt{* Expected Distance Sum: 12.57} \\
\texttt{* Penalties: 1.00} \\
\texttt{* Max Penalties: 8.00} \\
\texttt{* Rounds: 5.00} \\
\texttt{* Max Rounds: 12.00} \\
\texttt{* Object Count: 3.00} \\
            }
        }
    }
    & & \\ \\

\end{supertabular}
}

\end{document}
